
%=========== File containing some new commands ============%
%                                                          %
% Copyright (C) ISI - All Rights Reserved                  %
% Proprietary                                              %
% Written by Med Hossam <med.hossam@gmail.com>, April 2016 %
%                                                          %
% @author: HEDHILI Med Houssemeddine                       %
% @linkedin: http://tn.linkedin.com/in/medhossam           %
%==========================================================%

\newenvironment{changemargin}[2]{%
\begin{list}{}{%
\setlength{\leftmargin}{#1}%
\setlength{\rightmargin}{#2}%
}%
\item[]}
{\end{list}}

\makeatletter

%================= front cover variables =================%

\newcommand{\secondAuthor}[1]{\gdef\@secondAuthor{#1}}%
\newcommand{\@secondAuthor}{\@latex@warning@no@line{No \noexpand\secondAuthor given}}

\newcommand{\diplomaName}[1]{\gdef\@diplomaName{#1}}%
\newcommand{\@diplomaName}{\@latex@warning@no@line{No \noexpand\diplomaName given}}

\newcommand{\speciality}[1]{\gdef\@speciality{#1}}%
\newcommand{\@speciality}{\@latex@warning@no@line{No \noexpand\speciality given}}

\newcommand{\proFramerName}[1]{\gdef\@proFramerName{#1}}%
\newcommand{\@proFramerName}{\@latex@warning@no@line{No \noexpand\proFramerName given}}

\newcommand{\proFramerSpeciality}[1]{\gdef\@proFramerSpeciality{#1}}%
\newcommand{\@proFramerSpeciality}{\@latex@warning@no@line{No \noexpand\proFramerSpeciality given}}

\newcommand{\academicFramerName}[1]{\gdef\@academicFramerName{#1}}%
\newcommand{\@academicFramerName}{\@latex@warning@no@line{No \noexpand\academicFramerName given}}

\newcommand{\academicFramerSpeciality}[1]{\gdef\@academicFramerSpeciality{#1}}%
\newcommand{\@academicFramerSpeciality}{\@latex@warning@no@line{No \noexpand\academicFramerSpeciality given}}

\newcommand{\collegeYear}[1]{\gdef\@collegeYear{#1}}%
\newcommand{\@collegeYear}{\@latex@warning@no@line{No \noexpand\collegeYear given}}

\newcommand{\companyName}[1]{\gdef\@companyName{#1}}%
\newcommand{\@companyName}{\@latex@warning@no@line{No \noexpand\companyName given}}

%================== Signatures variables ==================%

\newcommand{\proSignSentence}[1]{\gdef\@proSignSentence{#1}}%
\newcommand{\@proSignSentence}{\@latex@warning@no@line{No \noexpand\proSignSentence given}}

\newcommand{\academicSignSentence}[1]{\gdef\@academicSignSentence{#1}}%
\newcommand{\@academicSignSentence}{\@latex@warning@no@line{No \noexpand\academicSignSentence given}}

%================== Backcover variables ==================%

\newcommand{\arabicAbstract}[1]{\gdef\@arabicAbstract{#1}}%
\newcommand{\@arabicAbstract}{\@latex@warning@no@line{No \noexpand\arabicAbstract given}}

\newcommand{\arabicAbstractKeywords}[1]{\gdef\@arabicAbstractKeywords{#1}}%
\newcommand{\@arabicAbstractKeywords}{\@latex@warning@no@line{No \noexpand\arabicAbstractKeywords given}}

\newcommand{\frenchAbstract}[1]{\gdef\@frenchAbstract{#1}}%
\newcommand{\@frenchAbstract}{\@latex@warning@no@line{No \noexpand\frenchAbstract given}}

\newcommand{\frenchAbstractKeywords}[1]{\gdef\@frenchAbstractKeywords{#1}}%
\newcommand{\@frenchAbstractKeywords}{\@latex@warning@no@line{No \noexpand\frenchAbstractKeywords given}}

\newcommand{\englishAbstract}[1]{\gdef\@englishAbstract{#1}}%
\newcommand{\@englishAbstract}{\@latex@warning@no@line{No \noexpand\englishAbstract given}}

\newcommand{\englishAbstractKeywords}[1]{\gdef\@englishAbstractKeywords{#1}}%
\newcommand{\@englishAbstractKeywords}{\@latex@warning@no@line{No \noexpand\englishAbstractKeywords given}}

\newcommand{\companyEmail}[1]{\gdef\@companyEmail{#1}}%
\newcommand{\@companyEmail}{\@latex@warning@no@line{No \noexpand\companyEmail given}}

\newcommand{\companyTel}[1]{\gdef\@companyTel{#1}}%
\newcommand{\@companyTel}{\@latex@warning@no@line{No \noexpand\companyTel given}}

\newcommand{\companyFax}[1]{\gdef\@companyFax{#1}}%
\newcommand{\@companyFax}{\@latex@warning@no@line{No \noexpand\companyFax given}}

\newcommand{\companyAddressFR}[1]{\gdef\@companyAddressFR{#1}}%
\newcommand{\@companyAddressFR}{\@latex@warning@no@line{No \noexpand\companyAddressFR given}}

\newcommand{\companyAddressAR}[1]{\gdef\@companyAddressAR{#1}}%
\newcommand{\@companyAddressAR}{\@latex@warning@no@line{No \noexpand\companyAddressAR given}}

%============= cmd for inserting blank page =============%
\newcommand\blankpage{%
    \null
    \thispagestyle{empty}%
    \addtocounter{page}{-1}%
    \newpage}

%================ document main language ================%
\selectlanguage{english}
%\selectlanguage{french}

%================== required packages ===================%

\usepackage{tcolorbox}
\usepackage{afterpage}
\usepackage{array,longtable,multirow}% http://ctan.org/pkg/{array,longtable,multirow}
\usepackage{pifont}

\usepackage{pdflscape}
\usepackage{rotating}
\usepackage{wrapfig}

%====== table packages used by chapters 1 & 2 ======%
\usepackage{booktabs}   % \toprule \midrule \bottomrule
\usepackage{tabularx}   % auto-width tabular* columns
\usepackage{makecell}   % \makecell multi-line cells

%====== code listings (inline snippets in chapters 3 & 4) ======%
% Clean monospaced listings for Dockerfiles, compose, CI workflow,
% Prisma schema and n8n JSON excerpts. The literate map remaps a few
% non-ASCII glyphs that pdfLaTeX cannot typeset verbatim.
\usepackage{listings}
\definecolor{codeBg}{HTML}{F5F5F5}
\lstset{
  basicstyle=\ttfamily\small,
  backgroundcolor=\color{codeBg},
  frame=single,
  rulecolor=\color{gray!45},
  breaklines=true,
  tabsize=2,
  showstringspaces=false,
  captionpos=b,
  extendedchars=true,
  literate=%
    {✅}{{\ding{52}}}1
    {❌}{{\ding{56}}}1
    {─}{{-}}1
    {—}{{-{}-}}2
    {…}{{...}}3
}

%====== realisation screenshot placeholder (chapters 3 & 4) ======%
% A sized, captioned empty frame standing in for an interface
% screenshot the author inserts later. Optional arg = box height
% (default 8cm); mandatory arg = description of the intended image.
% Built on tcolorbox (already loaded above) — no extra dependency.
\definecolor{phGray}{HTML}{F5F5F5}
\newcommand{\imagePlaceholder}[2][8cm]{%
  \begin{tcolorbox}[
      colback=phGray, colframe=gray!45, boxrule=0.5pt, arc=2pt,
      width=\linewidth, height=#1, halign=center, valign=center]
    \itshape\color{gray}[Screenshot to insert: #2]%
  \end{tcolorbox}%
}

%====== colours for the product backlog table (chapter 2) ======%
\definecolor{tableHeader}{RGB}{31,78,121}
% MoSCoW priorities
\definecolor{moscowMust}{HTML}{F8CBAD}
\definecolor{moscowShould}{HTML}{FFE699}
\definecolor{moscowCould}{HTML}{C6E0B4}
\definecolor{moscowWont}{HTML}{D9D9D9}
% complexity
\definecolor{complexLow}{HTML}{C6EFCE}
\definecolor{complexMed}{HTML}{FFEB9C}
\definecolor{complexHigh}{HTML}{FFD966}
\definecolor{complexVHigh}{HTML}{F4B183}
% feature bands
\definecolor{featAuth}{HTML}{DEEBF7}
\definecolor{featStore}{HTML}{E2EFDA}
\definecolor{featProduct}{HTML}{FCE4D6}
\definecolor{featOrder}{HTML}{FFF2CC}
\definecolor{featAdmin}{HTML}{EDEDED}
\definecolor{featAI}{HTML}{DDEDF8}
\definecolor{featPremium}{HTML}{E4DFEC}
\definecolor{featCICD}{HTML}{F2DCDB}
\definecolor{featAccount}{HTML}{E8E1F5}
\definecolor{featTheme}{HTML}{DCE3EC}

%====== clickable cross-references (hyperref) ======%
% Loaded last (after minitoc and caption). Makes every \ref, the table
% of contents / list of figures / list of tables, citations and URLs
% clickable links that jump to their target chapter, section, figure,
% table or reference. biblatex detects hyperref at end of preamble, so
% citation links work even though biblatex is loaded earlier.
\usepackage{hyperref}
\hypersetup{
  colorlinks=true,
  linkcolor=isiBlue,   % \ref to chapters, sections, figures, tables
  citecolor=isiBlue,   % \cite to the bibliography / webography
  urlcolor=isiBlue,    % \url and href
  linktoc=all,         % both the entry text and its page number link
  bookmarksnumbered=true,
  pdftitle={Design and Development of a Multi-Vendor E-Commerce Marketplace},
  pdfauthor={Nadim Jebali}
}

%========= English chrome: caption label override =========%
% The class hard-codes French caption labels (name = Tableau).
% With the English language selected above, override the table
% label to English. (Figures already use "Figure", identical in
% English.) Done here, after caption is loaded by the class.
% Table titles sit *below* the table (the \caption is emitted after the
% tabular in each table float); position=bottom applies the matching skip.
\captionsetup[table]{name=Table,position=bottom}